\section{Измерване на цената}
\label{sec:measurement}

Въпросът е колко струва разпознаването на една връзка в сравнение със сървър, който
предварително знае какво ще пристигне. Отговорът има две части. Измерването от край до
край показва дали цената се вижда отвън. Измерването в паметта показва колко е тя
всъщност.

\subsection{Постановка}

Сравняват се четири варианта на сървъра, всеки като отделен процес на една и съща машина:
\begin{itemize}
\item TLS с разпознаване: един сокет, разпознаване включено, клиентът говори TLS;
\item TLS с отделен слушател: разпознаването е изключено и сървърът започва
  ръкостискане веднага, без предварително четене;
\item HTTP с разпознаване: същото, но клиентът говори некриптиран HTTP;
\item HTTP с отделен слушател: разпознаването е изключено и сървърът подава връзката
  направо на анализатора на HTTP.
\end{itemize}

Всяка връзка носи една заявка и се затваря. Така цената за връзка се плаща веднъж на
заявка, а не се разпределя между много заявки по една дълга връзка. Клиентът е на Python
и използва стандартната проверка на сертификата и на името. Той отваря връзките една
след друга с темп \winRate{} в секунда. За всеки вариант и всяко повторение има по 20
загряващи връзки, които не влизат в резултата, и \winConns{} измерени. Повторенията са
\winReps{}, а редът на вариантите във всяко повторение се разбърква със записано начално
число (\winSeed{}), за да не попадне бавно изместване в състоянието на машината само в
един вариант.

Сървърът сам мери времената по монотонния си часовник и записва ред за всяка връзка.
Записват се четири момента: стартирането на нишката на връзката, решението за протокола,
готовността на транспорта и изпращането на отговора. Всички се отчитат от приемането на
връзката.

\subsection{Статистика}

Сравнява се медианата на времето до изпращане на отговора. Ефектът на разпознаването е
разликата на медианите, вариантът с разпознаване минус вариантът с отделен слушател.
Доверителният интервал от 95\% е получен с повторна извадка: 2000 пъти се теглят връзки
с връщане, поотделно в рамките на всяко повторение, и се взема разликата на медианите.
Интервалът е между 2,5-ия и 97,5-ия перцентил на тези разлики. Перцентилите на
разпределенията се изчисляват по метода на най-близкия ранг. Скриптът за анализ е в
хранилището и изходът му влиза в този документ автоматично, затова числата в текста не
са преписвани на ръка.

Освен интервала се показва и разликата между медианите на отделните повторения. Ако
повторенията се различават помежду си повече, отколкото е широк интервалът, интервалът
подценява неопределеността и това трябва да се каже.

\subsection{Първото измерване и какво показа то}

Първото измерване на Windows мереше времената от приемането на връзката. Разликите
излязоха $\winfirstTlsDiffA$~µs за TLS (интервал от $\winfirstTlsLowA$ до
$\winfirstTlsHighA$) и $\winfirstClearDiffA$~µs за HTTP (от $\winfirstClearLowA$ до
$\winfirstClearHighA$). Едно число в същите данни обаче показа, че тези разлики не
измерват разпознаването. При HTTP с отделен слушател медианата на времето до решението
беше \winfirstClearDedicatedDecideMedian{}~ms. Този вариант не чете нищо преди
решението, така че времето до него трябваше да е почти нула. Това, което се измерваше,
беше времето от приемането на връзката до момента, в който новата нишка започне да
работи.

Стартирането на нишка е еднакво за всички варианти и не е част от механизма. Затова във
второто измерване сървърът записва и момента на стартиране на нишката, и основният
показател става времето от него до изпращането на отговора. Първото измерване е запазено
в хранилището непроменено. То не е резултат, а измерването, което показа проблема.

\subsection{Резултати}

Таблици~\ref{tab:win} и~\ref{tab:lin} показват времето от стартирането на нишката до
изпращането на отговора, в милисекунди, за второто измерване на Windows и за Linux.

\begin{table}[H]
\centering
\caption{Windows, време от стартирането на нишката до отговора (ms)}
\label{tab:win}
\small
% Generated by bench/analyse.py. Do not edit by hand.
% metric: served_from_start_us, milliseconds
\begin{tabular}{lrrrrr}
\toprule
Рамо & $n$ & p50 & p90 & p99 & медиани по повторения \\
\midrule
TLS, с разпознаване & 2000 & 3{,}137 & 6{,}964 & 20{,}840 & 3{,}088\,\dots\,3{,}236 \\
TLS, отделен слушател & 2000 & 3{,}108 & 6{,}254 & 22{,}496 & 3{,}076\,\dots\,3{,}154 \\
HTTP, с разпознаване & 2000 & 0{,}126 & 0{,}436 & 20{,}424 & 0{,}124\,\dots\,0{,}128 \\
HTTP, отделен слушател & 2000 & 0{,}128 & 0{,}502 & 18{,}616 & 0{,}117\,\dots\,0{,}136 \\
\bottomrule
\end{tabular}

\end{table}

\begin{table}[H]
\centering
\caption{Linux във WSL2, време от стартирането на нишката до отговора (ms)}
\label{tab:lin}
\small
% Generated by bench/analyse.py. Do not edit by hand.
% metric: served_from_start_us, milliseconds
\begin{tabular}{lrrrrr}
\toprule
Рамо & $n$ & p50 & p90 & p99 & медиани по повторения \\
\midrule
TLS, с разпознаване & 2000 & 2{,}169 & 13{,}636 & 27{,}076 & 2{,}084\,\dots\,4{,}008 \\
TLS, отделен слушател & 2000 & 2{,}163 & 10{,}390 & 16{,}126 & 2{,}109\,\dots\,9{,}543 \\
HTTP, с разпознаване & 2000 & 0{,}084 & 0{,}116 & 0{,}242 & 0{,}072\,\dots\,0{,}095 \\
HTTP, отделен слушател & 2000 & 0{,}079 & 0{,}110 & 0{,}183 & 0{,}068\,\dots\,0{,}090 \\
\bottomrule
\end{tabular}

\end{table}

Ефектът на разпознаването, измерен от стартирането на нишката, е следният.
\begin{itemize}
\item Windows, TLS: $\winTlsDiffS$~µs, интервал от $\winTlsLowS$ до $\winTlsHighS$~µs,
  тоест $\winTlsRelS$\% от медианата на отделния слушател.
\item Windows, HTTP: $\winClearDiffS$~µs, интервал от $\winClearLowS$ до
  $\winClearHighS$~µs.
\item Linux, TLS: $\linTlsDiffS$~µs, интервал от $\linTlsLowS$ до $\linTlsHighS$~µs.
\item Linux, HTTP: $\linClearDiffS$~µs, интервал от $\linClearLowS$ до
  $\linClearHighS$~µs, тоест $\linClearRelS$\% от медиана от
  \linClearDedicatedMedian{}~ms.
\end{itemize}

Последната колона на таблиците показва най-малката и най-голямата медиана сред
отделните повторения. При TLS разликата между тях е по-голяма от самите ефекти, а при
HTTP е от същия порядък. При TLS на Linux интервалът на ефекта включва нулата.

\subsection{Цената в паметта}

Когато първото четене донася целия ClientHello, разпознаването не добавя системно
извикване. Остават заделяне на памет за буфера, копиране на ClientHello от буфера на
разпознаването в буфера на OpenSSL, сравнение на първия октет и създаване на
\code{ReplayReader}. Програмата \code{octet-micro} мери точно тази разлика, като
замества системното извикване с копиране от фиксиран буфер и в двата пътя. Така то се
съкращава по построение, а не чрез измерване. Резултатът е медиана от
\microWinRounds{} кръга по \microWinIterations{} повторения:
\begin{itemize}
\item Windows: \microWinDetecting{}~ns с разпознаване срещу \microWinDedicated{}~ns без,
  разлика \microWinDiff{}~ns на връзка;
\item Linux: \microLinDetecting{}~ns срещу \microLinDedicated{}~ns, разлика
  \microLinDiff{}~ns на връзка.
\end{itemize}

\subsection{Тълкуване}

Цената на разпознаването в паметта е под една микросекунда на връзка на двете
платформи. TLS връзката на същата машина отнема \winTlsDedicatedMedian{}~ms на Windows и
\linTlsDedicatedMedian{}~ms на Linux (медиана при отделен слушател), тоест цената е под
една хилядна от нея.

Измерването от край до край не може да различи цена от този порядък. Неговата
разделителна способност е десетки микросекунди, колкото е и разликата между
повторенията. Ефектите, които то показва, са в границите от няколко до около 30
микросекунди и не по-големи от 1\% при TLS.

Ефектът при HTTP на Linux, $\linClearDiffS$~µs, е с интервал, който не включва нулата,
и е с повече от порядък по-голям от цената в паметта. Причината за тази разлика не е
установена. Възможни обяснения са различният ред на изчакване и четене при двата
варианта или състоянието на кеша, но нито едно от тях не е проверено и работата не
приписва числото на разпознаването.

Измерването от приемането на връзката е изкривено от стартирането на нишките. Във
второто измерване на Windows медианите на това време бяха между
\winClearDedicatedStartMedian{} и \winTlsDetectStartMedian{}~ms в различните варианти, а
на Linux около \linClearDedicatedStartMedian{}~ms. Вариантите са изпълнявани в различни
моменти и затова са получили различно натоварване на машината. Разликите, измерени от
приемането, са $\winTlsDiffA$~µs при TLS и $\winClearDiffA$~µs при HTTP. Те измерват
машината, а не механизма, и не се цитират като цена на разпознаването.

Второто измерване на Windows отне \winDuration{}~s, а първото \winfirstDuration{}~s, при
еднакъв брой връзки. Причината за забавянето не е установена.

\subsection{Ограничения на измерването}

\begin{itemize}
\item Клиентът и сървърът са на една машина и говорят през локалния интерфейс. Мрежа
  няма.
\item Клиентът е на Python. Той е по-бавен от сървъра и времето, което сървърът чака
  неговите съобщения, влиза в измерените времена и на четирите варианта.
\item Linux е във WSL2, тоест във виртуална машина. Измерване на физически Linux няма.
\item Сървърът използва нишка за всяка връзка. Сървър с асинхронен вход и изход би имал
  други времена.
\item Първото измерване на Linux е отхвърлено: записите отиваха на диск на Windows през
  файловата система на WSL, която е бавна, и нишките на сървъра се натрупваха в очакване
  да запишат. То е запазено в хранилището под име, което казва това, и не е анализирано.
\end{itemize}
